\documentclass[11pt]{article}
\usepackage{amsmath,amssymb}
\usepackage[hidelinks]{hyperref}
\title{A supplier replacement can change the receive allowance without changing the receiver}
\author{}\date{}
\begin{document}\maketitle

\paragraph{What is being checked.}
The selected pinned SubT model has a separate rolling byte history at
each transmitter and receiver. In \texttt{attempt\_send}, however, both
comparisons use the capacity from the current \emph{sender's} radio
configuration, passed in by the broker. The policy parameter is not
the receive node's own configured capacity. This is a property of the
selected simulator source, not a claim about deployed radios. The
test reuses a named 173-byte MRWS completion and an independent 60-byte
survey packet addressed to the same station, plus separate traffic
to Home. Inputs include viable paths, deterministic draws and authority
to choose publication order before queue ownership passes to the broker.

\paragraph{Positive and negative cuts.}
A sender with threshold 200 and a sender with 233 can each send their
packet alone. With neither packet's old receive debit in range, both will
deliver exactly when the larger-cap sender's packet comes last. At cap
pair \((200,233)\) assigned to (completion, survey), completion first
works; at pair \((233,200)\), survey first works. The marginal receiver,
contents, total 233 bytes and sorted \emph{unassigned} capacities
are the same. A controller that only sees this view and must commit
once to an ordering cannot guarantee both receipts. A truthful
distinction of the two assignments before order commitment suffices,
or provision both possible smaller caps at 233, exactly. There is no
claim of obtaining this information in an untreated deployment.

An existing receive charge \(h\) adds to the required allowance when
it is still live; the replay spans the native inclusive age-one-second
boundary at ticks 4 and 5. Rejected receive leaves the sender bytes spent:
TX capacity was checked and charged before the RX comparison. The exact
application receipt, versus TX debit or a queue transition alone, is
the settlement boundary in this fixture. Changing the station's
configured cap among 1, 233 and 900 leaves the runs unchanged; changing
the sender's does not.

\paragraph{Reduction and law.}
Fix a dispatch cut; all mandatory packets are available, all their TX
attempts fit and reach their intended comparison, no bytes expire
during the cut and the order is freely selectable. Set \(p_i\) to size
and \(c_i\) to the qualified sender capacity. With initial load \(h\),
an order \(\pi\) succeeds on all receipts iff
\[
              h+\sum_{r\le j}p_{\pi_r}\le c_{\pi_j}
              \quad\text{for every }j .
\]
This is \href{https://ntrl.ntis.gov/NTRL/dashboard/searchResults/titleDetail/AD152722.xhtml}{Jackson's earliest-due-date rule}
with cumulative bytes as the processing coordinate and sender limits
as due dates. To see optimality, swap adjacent inverted caps \(c_a\ge c_b\)
in a feasible order: the maximum prefix through both was at most \(c_b\),
so the shorter new first prefix meets \(c_b\), and the full prefix still
meets \(c_a\). Later prefixes are unchanged. Repeat to sort. Sorting by
an abstract priority or by size alone does not replace these tests.

If a common order must work for a nonempty finite class of hidden
assignments, and sizes, initial receiver loads and destination sets are
fixed, substitute each row's minimum cap over that class. Taking a minimum
is exact here because every inequality must hold on every history and
its left side is already determined by the common order. It is not a
rule for uncertain sizes or receiver debit correlated with the cap, or
for choosing a sender after a new read. For multiple fixed receiver
sets, the same ascending-cap order solves all receiver subproblems
simultaneously when each packet has the same threshold at each of its
receivers. Each receiver sees a sorted subsequence, and each subsequence
must pass its own initial-debit check. The generalization assumes the
total per-sender TX attempt bills fit despite fanout; receiver-specific
cap order or forced precedence can invalidate the reduction.

\paragraph{Evidence and limits.}
1,296 pinned callback-and-debit cases match the direct sender-threshold
formula, including unsuccessful receipt with charged TX. Independently
enumerating all orders agrees in 520 each of one-receiver, paired-world
and fixed multi-receiver model cases. Lean checks the adjacent exchange,
the paired obstruction, the sharp two-world capacity threshold and
pointwise minimum law. The strongest applicable ordering baseline is
the same rule, not a new scheduler. Packet creation uses actual selected
warehouse completion; the physical or operational appointment, cap
reporting and order authority are not instantiated here. In-block
expiry, multiple retries, reset/membership and pose blindness each need
separate premises. No independent mission or programme acceptance.
\end{document}
